\section{Requirement coverage and report continuity}
Table~\ref{tab:coverage} maps the M3 expectations in the CSE 516 project
guidelines to the evidence in this report and repository.

\begin{table}[!htbp]\centering\small
\begin{tabularx}{\linewidth}{@{}lX@{}}\toprule
M3 expectation & Evidence in the cumulative report/package\\\midrule
Full experimental results & Five main seeds; three stress-test seeds; raw metrics and configs\\
Baseline comparison & M2 means/SGD/ADF/EP; graph-free MAP and matching GMRF\\
KPI plots and tables & Rating/probability errors, cost, reliability, convergence, stress and errors\\
Variation and uncertainty & Sample SD; 2,000-replicate paired user-bootstrap intervals with limits\\
Probability quality & NLL, Brier, ECE, reliability and interval outputs\\
Robustness & Training availability, label corruption, assumed noise, graph density, cold items/users\\
Ablation and error analysis & Gamma=0, genre/residual components, rating/support groups\\
Interpretation & Main-result explanations, counterexamples and explicit attribution limits\\
Theoretical validation & Preserved EP derivation; proper GMRF prior; exact/Gibbs/iterative diagnostic; quadrature\\
Revised SOTA position & Specific Matchbox/PreRec models and protocols; no direct SOTA score claim\\
Reproducibility & Source/config/environment hashes, run logs, tests, result/figure scripts\\
Team roles & Member responsibilities listed in the team section below\\\bottomrule
\end{tabularx}
\caption{Coverage of the M3 guidelines. SHD is inapplicable without a true graph.}\label{tab:coverage}
\end{table}

The updated source begins from the revised M2 \code{main.tex} and preserves its
title, team details, model equations, EP/VMP derivation and numerical-validation
discussion. It updates the abstract, milestone/date, benchmark, results,
positioning and conclusion; it adds the GMRF representation and M3 analysis.
A detailed M2-to-M3 changelog is in \code{Review/CHANGELOG\_M2\_TO\_M3.md}.

Not every item mentioned as possible future work in M2 is claimed complete.
The implemented additions cover noise/prior/damping/order sensitivity, longer
EP runs, genre ablations, reduced data and cold-start evaluations. Learned
user-specific thresholds, a temporal deployment test, and full-posterior GMRF
inference remain deferred. The formal general guidelines encourage depth in
one question, rather than an unsupported claim to have tested every PGM method.

\section{Team accountability and AI assistance}
Team roles follow the M1/M2 plan: Jyotier Vithalani (AU2440170): mathematical
model, EP/ADF and repository; Tanmay Kanani (AU2440234): baselines; Ansh
Chaudhari (AU2440259): paper/SOTA analysis; Shlok Shah (AU2440178): dataset and
metadata; Het Patel (AU2440221): KPIs, experiment tracking and documentation.

Generative AI assisted implementation, execution, analysis and writing. All
numerical results were computed by the included code on the verified dataset;
team members remain responsible for checking them.

\section{Reproduction and evidence boundaries}
Code, configurations and results are in the
\code{M3\_g4\_th\_apl\_recommend-sys\_team} folder of the project repository.
The Code README provides commands for validation, main/stress/sensitivity
experiments, figure regeneration and report compilation. Raw MovieLens files
are not redistributed; the official downloader, source links and hashes are
included. Seed-42 checkpoints and ID mappings support direct probability queries.

The main-run source snapshot and its hashes preserve the exact executed
implementation. Later changes add a cold-user condition and audit metadata to
the robustness path, without changing the main numerical kernels or selection
rules. Results from the separate main, robustness and sensitivity stages retain
their individual configurations and source/environment records.

The earlier synthetic offline demo is excluded from all MovieLens benchmark
tables. The 12-variable Gaussian example is labelled a diagnostic throughout.
Historical M2 files are labelled separately. Training time excludes JIT and
test evaluation; comparison of complete workflows additionally includes graph
construction and validation search. Test evaluation follows validation selection
for each comparison.
